\documentclass[10pt,a4paper]{amsart}
\usepackage[margin=19mm]{geometry}
\usepackage{amsmath,amssymb,mathtools}
\usepackage[scr=rsfs,cal=euler]{mathalfa}
\usepackage{xfrac}
\usepackage[unicode,hidelinks,bookmarks=false]{hyperref}
\newcommand{\ie}{i.e.\ }
\newcommand{\cf}{cf.\ }
\newcommand{\resp}{respectively\ }
\newcommand{\Subsection}{\subsection}
\providecommand{\nobreakdash}{\nobreak\hbox{-}}
\setlength{\emergencystretch}{2em}
\multlinegap0pt
\usepackage{bm}
\usepackage{bbm}
\usepackage{dsfont}
\usepackage{xspace}
\usepackage{cancel}
\usepackage{mathtools}
\usepackage{cleveref}
\usepackage{amsthm}
\makeatletter
\@ifundefined{definition}{\newtheorem{definition}{Definition}}{}
\@ifundefined{lemma}{\newtheorem{lemma}{Lemma}}{}
\@ifundefined{theorem}{\newtheorem{theorem}{Theorem}}{}
\makeatother
\title[A family of accumulation points of non-free rational numbers]{A family of accumulation points of non-free rational numbers}
\author{Christopher Buyalos, Jayden Thadani, Xinbei Wang, Bradley Zykoski, Michael Zshornack}
\date{Source: arXiv:2511.10422v2 (CC BY 4.0)}
\begin{document}
\maketitle

\noindent\emph{Source and licence.} A family of accumulation points of non-free rational numbers. Version arXiv:2511.10422v2 (CC BY 4.0), distributed under the Creative Commons Attribution 4.0 International licence, as recorded in the arXiv licence metadata for this exact version and shown on the abstract page https://arxiv.org/abs/2511.10422v2. Full rights notice: licenses/arxiv-2511.10422-CC-BY-4.0.txt.\par\smallskip

\noindent\textit{Editorial scope.} Counted unit: the Theorem whose source label is \texttt{thm:main} in the article (source0.tex lines 210--213, source label \texttt{thm:main}), delivered together with the article's own complete original proof of it (source0.tex lines 240--269, one \texttt{proof} environment of 30 lines). No mathematics is written, repaired or re-derived: statement and proof are the article's own text line for line. The proof is detached from the statement in the source: the article states the result at lines 210--213 and prints its proof at lines 240--269, and the proof environment's own header names the statement, so the association is the article's own. Required context retained uncounted: def:curve (lines 219--230); lem:dioph (lines 234--236). Every definition, display and label of those lines is retained unchanged. Omitted from this package and not redistributed: 1--209, 214--218, 231--233, 237--239, 270--450. Nothing inside a retained excerpt is omitted or altered: every retained body is delivered exactly as the article's lines are written, so no omission receipt of any kind accompanies this package. Citations: the retained excerpts carry no citation command at all, so no bibliography entry of the article is transcribed and no reference is redistributed. Reference adaptation: none was needed and none was made; every reference inside the delivered unit resolves inside it, so no unresolved reference or citation is produced. Printed slips kept literal: no printed word, symbol, hypothesis or number of the counted statement or of its proof was altered. Layout and class: the article's own class is \texttt{amsart} with packages including inputenc, article-preamble, todonotes; the wrapper is the lane's pinned \texttt{amsart} editorial wrapper at 10pt on A4, which restates the theorem-like environments the retained text needs and adds the article's own macro definitions that the retained text needs (none), with extra wrapper packages (bm, bbm, dsfont, xspace, cancel, mathtools, cleveref). The delivered numbering is the unit's own. Assets: no retained span contains a figure environment, a graphics-inclusion command, a PDF-page inclusion, a table or a TikZ picture; no image file is copied into this package, the package declares no asset dependency and no image file is redistributed with it. Licence: version arXiv:2511.10422v2 of this article is distributed under the Creative Commons Attribution 4.0 International licence, as recorded in the arXiv licence metadata for this exact version and shown on the abstract page https://arxiv.org/abs/2511.10422v2. A full rights notice is delivered at licenses/arxiv-2511.10422-CC-BY-4.0.txt. Characters: every retained excerpt is pure ASCII, so the delivered engine, XeLaTeX with its default Latin Modern text font, reports no missing character.
\begin{definition}
\label{def:curve}
Let $a = (a_{1}, \dots, a_{5}) \in \mathbb{Z}^{5}$. We define a curve $C_a$ as follows. The discriminant of $P_{\text{HR}}^{5}(a_{1}, \dots, a_{5}; q)$ is
\[
	\Delta_{a} = (a_{3} a_{4} a_{5})^{2} + (2 a_{3} a_{4}^{2} a_{5}^{2} - 2 a_{2} a_{3} a_{4} a_{5}^{2} - 2 a_{2} a_{3}^{2} a_{4} a_{5}) a_{1} + ((a_{2} a_{3}^{2}) + (a_{2} a_{5})^{2} + 2 a_{2} a_{4} a_{5}^{2} + 2 a_{2}^{2} a_{3} a_{5} - 2 a_{2} a_{3} a_{4} a_{5}) a_{1}^{2}
\]
If we consider $a_2, a_3, a_4, a_5$ to be fixed and $a_1$ to be variable, then the discriminant $\Delta_{a}$ is a square integer if and only if there exists an integer $y$ such that $(x, y) = (a_1, y)$ is a root of a polynomial of the form
\[
	F_{a}(x, y) = \alpha_{a} x^{2} + \beta_{a} x + \gamma_{a} - y^{2}.
\]
We define $C_{a}$ to be the zero locus of $F_{a}$.
\end{definition}

\begin{lemma} \label{lem:dioph}
	For each non-zero  $a_{3}, a_{4}, a_{5} \in \mathbb{Z}$ there are infinitely many $a_{1}, a_{2} \in \mathbb{Z}$ such that $C_{a}$ has infinitely many integer points.
\end{lemma}

\begin{theorem}
\label{thm:main}
All $1$-step relation numbers are accumulations of non-constant sequences of length-$5$ rational half-relation numbers.
\end{theorem}

\begin{proof}[Proof of \cref{thm:main}] 
Let $a = (a_{1}, \dots, a_{5}) \in \mathbb{Z}^{5}$. The length $5$ half-relation numbers are exactly the roots $q_{a}^{+}, q_{a}^{-}$ of the polynomials
\begin{align*}
	P_{\text{HR}}^5(a_1, a_2, a_3, a_4, a_5; q) & =  (a_1a_2a_3a_4a_5)q^2+(a_1a_2a_3-a_2a_3a_4+a_1a_2a_5+a_1a_4a_5+a_3a_4a_5)q\\
     & +(a_1-a_2+a_3-a_4+a_5) \\
     & = c_{2} q^{2} + c_{1} q + c_{0} 
\end{align*}
for $a \in \mathbb{Z}^{5}$. Specifically, let
\begin{align*}
	q_{a}^{+} & = \frac{- c_{1} + \sqrt{c_{1}^{2} - 4 c_{0} c_{2}}}{2 c_{0}} & q_{a}^{-} & = \frac{- c_{1} - \sqrt{c_{1}^{2} - 4 c_{0} c_{2}}}{2 c_{0}}
\end{align*}
The roots $q_{a}^{\pm}$ are rational exactly when the discriminant of $P_{\text{HR}}^{5}(a_{1}, \dots, a_{5}; q)$ is a square integer. 

\cref{def:curve} expresses the condition that $\Delta_{a}$ is a square integer as an integer point on the curve $C_{a}$. In \cref{lem:dioph}, we show that for each $a_{3}, a_{4}, a_{5} \in \mathbb{Z}$, there are infinitely many $a_{1}, a_{2} \in \mathbb{Z}$ such that $C_{a}$ has infinitely many integer points. We claim that this suffices to prove \cref{thm:main}. 

For each $a_{3}, a_{4}, a_{5}$, there are sufficiently many (infinitely many) rational roots $q_{a}^{-}$ of $P_{\text{HR}}^{5}(a_{1}, \dots, a_{5}; q)$. Then the limit of the rational $q_{a}^{-}$ agrees with the limit of all $q_{a}^{-}$. Specifically, each
\begin{align*}
	\lim_{a_1\rightarrow \infty} q_a^{\pm} & = \frac{-\big(a_2a_3+a_2a_5+a_4a_5\big) \pm \sqrt{(a_2a_3)^2+(a_2a_5)^2+(a_4a_5)^2+2a_2a_4a_5^2+2a_2^2a_3a_5-2a_2a_3a_4a_5}}{2a_2a_3a_4a_5}
\end{align*}
and each
\begin{align*}
	\lim_{a_{2} \to \infty} \lim_{a_{1} \to \infty} q_{a}^{-} & = \frac{-(a_{3} + a_{5})}{a_{3} a_{4} a_{5}} & \lim_{a_{2} \to \infty} \lim_{a_{1} \to \infty} q_{a}^{+} & = 0
\end{align*}
is a limit of rational length $5$ half relation numbers. Note that even on taking further limits, we only get $1$-step relation numbers as limits --- we have
\[
	\lim_{a_{3} \to \infty} \lim_{a_{2} \to \infty} \lim_{a_{1} \to \infty} q_{a}^{-} =  \frac{-1}{a_{4} a_{5}}
\]

Here, we chose a specific order in which to take limits. However, this order is not necessary to obtain the $1$-step relation numbers of Kim--Koberda as limit points.
\end{proof}

\end{document}
